\documentclass[11pt,a4paper]{letter}
\usepackage[T1]{fontenc}
\usepackage[utf8]{inputenc}
\usepackage[margin=2.5cm]{geometry}
\usepackage{microtype}

\signature{Arun V.\ Eluvathingal\\
           Department of Electrical Engineering\\
           IIT Madras, Chennai 600036, India\\
           \texttt{arun.eluvathingal@iitm.ac.in}}

\address{Arun V.\ Eluvathingal\\
         Department of Electrical Engineering\\
         IIT Madras, Chennai 600036, India}

\begin{document}
\begin{letter}{Editor-in-Chief\\
               \textit{IET Generation, Transmission \& Distribution}}

\opening{Dear Editor,}

We submit for your consideration the manuscript:

\begin{center}
  \textbf{Adaptive Forgetting Factor for Drifting Inverter-Based Resource
  Penetration in Model-Based Protection Systems}
\end{center}

\noindent for publication in \textit{IET Generation, Transmission \& Distribution}.

\medskip
\noindent\textbf{Context.}
The Self-Adaptive Model-Based Protection System (SAMBPS) uses a
Levenberg--Marquardt zone-model estimator to adapt relay settings to the
instantaneous IBR current ratio $k_{\rm ibr}$.  Existing SAMBPS publications
assume $k_{\rm ibr}$ is approximately stationary within the estimation window.
In practice, solar PV follows a diurnal ramp, wind generation executes rapid
fluctuations, and topology switching creates step discontinuities — causing
$k_{\rm ibr}$ to drift on timescales of seconds to hours.  A fixed forgetting
factor ($\lambda = 0.995$) creates a 63\,ms tracking lag that degrades protection
detection probability by up to 16\% during drift.

\medskip
\noindent\textbf{Key contributions.}
\begin{enumerate}
  \item \textbf{CUSUM drift detector.}
        A one-sided CUSUM applied to successive LM estimates $\hat{k}_{\rm ibr}$
        detects drift within 3 cycles (median) for all four drift archetypes,
        with false-alarm rate $< 2\!\times\!10^{-4}$ per cycle under
        noise-only conditions.
  \item \textbf{Logistic $\lambda$-adaptation law.}
        A closed-form law interpolates between $\lambda_{\min} = 0.970$
        (effective window 33 samples, fast tracking during drift) and
        $\lambda_{\max} = 0.995$ (200 samples, low noise during stationary
        conditions), driven by the estimated drift rate $\dot{r}_k$.
  \item \textbf{Stability proposition.}
        A formal proof that $\kappa_n(\mathbf{J}) < 30$ is preserved for all
        $\lambda_k \geq \lambda_{\min} = 0.970$, guaranteeing LM convergence and
        the SAMBPS Stage-2 veto integrity under all credible drift conditions.
  \item \textbf{Monte Carlo validation.}
        20\,000 trials across four drift scenarios (solar diurnal, cloud transient,
        wind random walk, topology step) confirm: tracking RMSE reduced $3.4\times$
        (0.062 to 0.018\,pu); $P_{\rm FA}$ reduced to $\leq 0.002$;
        $P_D \geq 0.981$ vs.\ $0.789$--$0.862$ for fixed $\lambda$.
\end{enumerate}

\medskip
\noindent\textbf{Relationship to companion papers.}
This is gap paper G2 in the SAMBPS series.  The SAMBPS framework is developed
in companion papers submitted to \textit{IEEE Transactions on Smart Grid}
(unified architecture, J4) and \textit{IEEE Transactions on Power Delivery}
(87L algorithm, J1; SEL-411L hardware validation, G1).  This paper addresses
a distinct open problem — the non-stationary $k_{\rm ibr}$ environment — that
arises between fault events and is not covered in the fault-scenario validation
of the companion papers.

\medskip
\noindent\textbf{Suggested reviewers.}
\begin{itemize}
  \item Prof.\ Ali Hooshyar, University of Toronto
        (IBR protection, adaptive relay settings)
  \item Prof.\ Wilsun Xu, University of Alberta
        (power quality, IBR system identification)
  \item Prof.\ Geza Joos, McGill University
        (power electronics, inverter-grid interaction)
  \item Prof.\ Lennart Ljung, Linköping University
        (system identification, recursive estimation)
\end{itemize}

\medskip
\noindent
The manuscript has not been submitted elsewhere and all authors have approved
the final version.

\closing{Sincerely,}
\end{letter}
\end{document}
